\documentclass{report}
\usepackage{amsmath,amssymb}
\setlength{\parindent}{0mm}
\setlength{\parskip}{1em}
\begin{document}
\begin{center}
\rule{6in}{1pt} \
{\large Dmitriy Anistratov \\
{\bf Multilevel Iteration Methods for Eigenvalue Neutron Transport Problems}}

2105 Burlington Labs \\ Dept of Nuclear Engineering \\ NC State University \\ Raleigh \\ NC 27695-7909
\\
{\tt anistratov@ncsu.edu}\end{center}

In a variety of nuclear reactor design calculations, one needs to
evaluate a critical parameter of a reactor and the corresponding particle
distribution function. To perform computations for this important class
of reactor physics problems, it is necessary to solve the neutron
transport equation and determine the largest eigenvalue with the
associated eigenfunction. The transport equation is a detailed
conservation relation for particles in the phase space defined by the
spatial, angular and energy variables. It is a linear
integro-differential equation that can be presented in the following
form:
\[
H\psi=S_s\psi+ \frac{1}{k}S_f\psi \, , \quad \psi=\psi(\vec r, \vec \Omega, E) \, ,
\]
Here the operator $H$ accounts for particle streaming and collisions. It
is a differential operator of the first order that is not a self-adjoint
one. The operators $S_s$ and $S_f$ are integral operators that describe
scattering and fission processes. $k$ is the eigenvalue.


To solve eignvalue transport problems, we use nonlinear multilevel
iterative methods. The basic idea behind these methods is to effectively
reduce the dimensionality of the problem by averaging the transport
equation over angular and energy variables. We apply two successive
projection operations to the transport equation with no approximations
involved. This procedure leads to a multilevel system of equations in
spaces with lower dimensionality than the original space to which the
transport solution belongs. The multilevel system of equations is closed
by means of exact closure relations using linear-fractional factors. The
solution of transport equation provides all necessary data that enable
one to define the transport problems of the reduced dimensionality.


In the presented method, the eigenvalue is evaluated in the subspace of
the lowest dimensionality where solution depends only on the spatial
variable. Note that this is an exact projected solution subspace, because
all closure relations defined to derive the multilevel equations are
exact. The resulting low-order transport problem is treated as the
generalized eigenvalue problem. This low-order transport problem
reproduces essential features of neutron transport physics. As a result,
one gets a good estimation of the eigenvalue and associated
eigenfunction. This leads to acceleration of the multigroup transport
iterations.

We apply this methodology to the nonlinear diffusion acceleration (NDA)
method that is an efficient and flexible transport iterative scheme for
solving reactor-physics problems. In this talk we present a fast
iterative algorithm for solving multigroup neutron transport eigenvalue
problems in 1D slab geometry. Numerical results that illustrate
performance of the algorithm are demonstrated.


\end{document}
